\notepage{11}{Hermitian operators: eigenvalues and orthogonality}
\smallheading{Degeneracy}
A repeated root of the characteristic equation
$\det(\Omega-\omega I)=0$ has an \emph{algebraic multiplicity}.
A double or triple root has multiplicity 2 or 3. In the Hermitian case,
this equals the dimension of the corresponding eigenspace, called
its degeneracy.

If $\ket{v_1},\ldots,\ket{v_m}$ are independent eigenvectors with the
same eigenvalue $\omega$, then any nonzero combination satisfies
\[
 \Omega\sum_{a=1}^{m}a_a\ket{v_a}
 =\omega\sum_{a=1}^{m}a_a\ket{v_a}.
\]
Thus every nonzero vector in their span is an eigenvector with
eigenvalue $\omega$.

\smallheading{Theorem 10}
Every finite-dimensional Hermitian operator has real eigenvalues and
admits a complete orthonormal set of eigenvectors.

To prove reality and orthogonality, suppose
\[
 \Omega\ket{\omega_i}=\omega_i\ket{\omega_i},\qquad
 \Omega\ket{\omega_j}=\omega_j\ket{\omega_j}.
\]
These equations also allow $i=j$, or $i\ne j$ with
$\omega_i=\omega_j$. Taking matrix elements gives
\begin{align*}
 \mel{\omega_j}{\Omega}{\omega_i}
 &=\omega_i\braket{\omega_j}{\omega_i},\\
 \mel{\omega_i}{\Omega}{\omega_j}
 &=\omega_j\braket{\omega_i}{\omega_j}.
\end{align*}
Conjugating the second equation and using
$\Omega^\dagger=\Omega$ yields
\begin{align*}
 \mel{\omega_j}{\Omega}{\omega_i}
 &=\omega_j^*\braket{\omega_j}{\omega_i}.
\end{align*}
Subtracting the two expressions,
\[
 \boxed{(\omega_i-\omega_j^*)\braket{\omega_j}{\omega_i}=0.}
\]
The consequences for equal and unequal eigenvalues follow on note~12.
